\section{Step 86 Select compact codes when available}\label{h263-coeff-s06}
Look up LAST, RUN and the positive magnitude in the supplied complete table.
When found, use its code and declared width, setting the final sign bit for a
negative level. Otherwise retain the escape record. Keep both implementations
for correctness and later size/work comparison.

A code's integer value does not encode its width. The compact final $+1$ code
has leading zero, so passing only its nonzero bit length would change the syntax.

\paragraph{Observe} The default compact block uses sixteen bits; the escape-only
block uses fifty-two. Both describe the same DC and AC levels.
\Needspace{8\baselineskip}
\paragraph{Check} What identifies one compact table entry?
\begin{itemize}
\item[A.] Only the sign of the level.
\item[B.] The current raw pixel value.
\item[C.] LAST, RUN and level magnitude.
\end{itemize}
